\section{Conclusions}
\label{sec:Conclusions}

In this paper, we presented a modular, domain-agnostic open-source software architecture that connects conversational Natural Language Understanding (NLU), Apache Solr spatial indexing, and dynamic Geographic Information System (GIS) map visual synchronization. The architecture eliminates database schema hallucinations by coupling Few-Shot intent parsing with strict OpenAPI JSON Schema constraints evaluated under greedy decoding. 

The software translates conversational spatial queries into structured Solr Boolean filter parameters, synchronizing visual filter badges and animated pulse markers across an interactive Leaflet web map. Validated within the European research project \textbf{CRMsDataSpace} as a demanding strategic domain case study, empirical evaluation across 40 critical stress tests and 100 benchmark queries on an institutional NVIDIA A100 testbed demonstrated high extraction fidelity (100.0\% intent accuracy, 100.0\% country and commodity F1-scores, and up to 93.4\% overall macro F1-score across open-weight models including Qwen~2.5 7B, Llama~3.2 3B, and DeepSeek~R1 7B). By enabling native on-premises local execution under Slurm without external cloud API dependencies, the architecture guarantees 100\% European data sovereignty in compliance with Regulation (EU) 2024/1252 and GAIA-X principles, while the zero-setup standalone mock execution mode and synthetic 100-site dataset ensure immediate out-of-the-box reproducibility for academic reviewers and practitioners.

Future developments will focus on expanding the architecture toward multi-agent coordination, integrating Open Subsurface Data Universe (OSDU) standards for deep analytical modeling, and supporting dynamic raster layer overlays (GeoTIFF) directly within the conversational GIS feedback loop.
